
\subsubsection{Semantic Uncertainty Estimation}

Farquhar et al. (2024) introduced semantic entropy for detecting hallucinations in large language models. The method generates K samples at temperature T, clusters semantically equivalent responses using natural language inference, and computes entropy over meaning clusters rather than token sequences. Validation on TruthfulQA and other benchmarks established AUROC exceeding 0.80. The approach requires approximately 1500ms per query due to multi-sample generation followed by DeBERTa-based clustering.

Kossen et al. (2024) proposed training lightweight probes on hidden states to predict semantic entropy, achieving AUROC approximately 0.80 while avoiding multi-sample generation during inference. The method requires labeled training data (semantic entropy values themselves require expensive computation for the training set) and produces black-box predictions.

Manakul et al. (2023) developed SelfCheckGPT, which measures self-consistency across 5-20 samples using BERTScore, n-gram overlap, or NLI comparison. The method achieves AUROC approximately 0.75 on question-answering datasets but incurs 5-20 times latency overhead. Conformal prediction approaches (Kumar et al., 2023) provide statistical guarantees but similarly require multi-sample generation.

\subsubsection{Spectral Analysis of Neural Networks}

The NerVE framework applied participation ratio and eigenvalue spectrum analysis to feed-forward network weights, finding that geometric properties correlate with model behavior (ICLR 2026). However, NerVE analyzes static weight geometry rather than per-example hidden state dynamics. Voita et al. (2019) demonstrated that transformer representations contain redundant dimensions, with effective dimensionality varying across layers. Ashukha et al. (2020) established benchmarks for ensemble-based uncertainty estimation.

\subsubsection{Computational Efficiency}

Token-level perplexity provides fast uncertainty estimates but lacks semantic understanding. Rahmati et al. (2025) proposed C-LoRA for contextual uncertainty estimation through low-rank adaptation, though this requires supervised training. Wang et al. (2025) introduced GENUINE, achieving 29\% higher AUROC than semantic entropy with 15\% lower calibration error, though computational costs for 7B+ models remain uncharacterized.

Our work hypothesized that intrinsic geometric properties computable from single forward passes could provide training-free uncertainty proxies. The investigation revealed that hidden state extraction itself creates a computational bottleneck.
